%%%%%%%%%%%%%%%%%%%%%%%%%%%%%%%%%%%%%%%%%%%%%%%%%%%%%%%%%%%%%%%%%%%%%%%%%%%%%%%%%%%%%%%%%%%%%%%%%%%%%%%%%%%%%%%%%%%%%%%%%%%%%%%%%%%%%%%%%%%%%%%%%%%%%%%%%%%
% This is just an example/guide for you to refer to when submitting manuscripts to Frontiers, it is not mandatory to use Frontiers .cls files nor frontiers.tex  %
% This will only generate the Manuscript, the final article will be typeset by Frontiers after acceptance.                                                 %
%                                                                                                                                                         %
% When submitting your files, remember to upload this *tex file, the pdf generated with it, the *bib file (if bibliography is not within the *tex) and all the figures.
%%%%%%%%%%%%%%%%%%%%%%%%%%%%%%%%%%%%%%%%%%%%%%%%%%%%%%%%%%%%%%%%%%%%%%%%%%%%%%%%%%%%%%%%%%%%%%%%%%%%%%%%%%%%%%%%%%%%%%%%%%%%%%%%%%%%%%%%%%%%%%%%%%%%%%%%%%%

%%% Version 3.1 Generated 2015/22/05 %%%
%%% You will need to have the following packages installed: datetime, fmtcount, etoolbox, fcprefix, which are normally inlcuded in WinEdt. %%%
%%% In http://www.ctan.org/ you can find the packages and how to install them, if necessary. %%%

\documentclass{frontiersSCNS} % for Science, Engineering and Humanities and Social Sciences articles
%\documentclass{frontiersHLTH} % for Health articles
%\documentclass{frontiersFPHY} % for Physics and Applied Mathematics and Statistics articles

%\setcitestyle{square}
\usepackage{url,hyperref,lineno,microtype}
\usepackage[onehalfspacing]{setspace}
\linenumbers


% Leave a blank line between paragraphs instead of using \\


\def\keyFont{\fontsize{8}{11}\helveticabold }
\def\firstAuthorLast{Stewart {et~al.}} %use et al only if is more than 1 author
\def\Authors{Terrence C. Stewart\,$^{1,*}$, Ashley Kleinhans\,$^{2}$ and Chris Eliasmith\,$^1$}
% Affiliations should be keyed to the author's name with superscript numbers and be listed as follows: Laboratory, Institute, Department, Organization, City, State abbreviation (USA, Canada, Australia), and Country (without detailed address information such as city zip codes or street names).
% If one of the authors has a change of address, list the new address below the correspondence details using a superscript symbol and use the same symbol to indicate the author in the author list.
\def\Address{$^{1}$Centre for Theoretical Neuroscience, University of Waterloo,
    Waterloo, ON, Canada
    \\
$^{2}$Mobile Intelligent Autonomous Systems group, 
      Council for Scientific and Industrial Research,
      Pretoria, South Africa}
% The Corresponding Author should be marked with an asterisk
% Provide the exact contact address (this time including street name and city zip code) and email of the corresponding author
\def\corrAuthor{Terrence C. Stewart}
\def\corrAddress{Centre for Theoretical Neuroscience, University of Waterloo,
    Waterloo, ON, Canada}
\def\corrEmail{tcstewar@uwaterloo.ca}




\begin{document}
\onecolumn
\firstpage{1}

\title[Embodied Neuromorphic Benchmarks]{Embodied Neuromorphic Benchmarks} 

\author[\firstAuthorLast ]{\Authors} %This field will be automatically populated
\address{} %This field will be automatically populated
\correspondance{} %This field will be automatically populated

\extraAuth{}% If there are more than 1 corresponding author, comment this line and uncomment the next one.
%\extraAuth{corresponding Author2 \\ Laboratory X2, Institute X2, Department X2, Organization X2, Street X2, City X2 , State XX2 (only USA, Canada and Australia), Zip Code2, X2 Country X2, email2@uni2.edu}

\maketitle

%%%%%%%%%%%%%%%%%%%%%%%%%%%%%%%%%%%%%%%%%%%%%%%%%%%%%%%%%%%%%%%%%%%%%%%%%%%%%%%%%%%%%%%%%%%%%%%%%%%%%%%%%%%%%%%%%%%%%%%%%%%%%%%%%%%%%%%%%%%%%%%%%%%%%%%%%%%%%%%%%%%%%%%%%%%%%%%%%%%%%%%%%%%%%%%%%%%%%%%%%%%%%%%%%%%%%%%%%%%%%%%%%%%%%%%
%%% The sections below are for reference only.
%%%
%%% For Original Research Articles, Clinical Trial Articles, and Technology Reports the section headings should be those appropriate for your field and the research itself. It is recommended to organize your manuscript in the
%%% following sections or their equivalents for your field:
%%% Abstract, Introduction, Material and Methods, Results, and Discussion.
%%% Please note that the Material and Methods section can be placed in any of the following ways: before Results, before Discussion or after Discussion.
%%%
%%%For information about Clinical Trial Registration, please go to http://www.frontiersin.org/about/AuthorGuidelines#ClinicalTrialRegistration
%%%
%%% For Clinical Case Studies the following sections are mandatory: Abstract, Introduction, Background, Discussion, and Concluding Remarks.
%%%
%%% For all other article types there are no mandatory sections.
%%%%%%%%%%%%%%%%%%%%%%%%%%%%%%%%%%%%%%%%%%%%%%%%%%%%%%%%%%%%%%%%%%%%%%%%%%%%%%%%%%%%%%%%%%%%%%%%%%%%%%%%%%%%%%%%%%%%%%%%%%%%%%%%%%%%%%%%%%%%%%%%%%%%%%%%%%%%%%%%%%%%%%%%%%%%%%%%%%%%%%%%%%%%%%%%%%%%%%%%%%%%%%%%%%%%%%%%%%%%%%%%%%%%%%%

\begin{abstract}

    Evaluating the effectiveness and performance of neuromorphic hardware is
    difficult.  It is even more difficult when the task of interest is an
    embodied task; that is, a task where the output from the neuromorphic
    hardware affects its future input.  However, these embodied situations
    are one of the primary potential uses of neuromorphic hardware.  To address
    this, we present a methodology for embodied benchmarking and a particular
    example of a benchmark for adaptive control of an arbitrary system.  These
    benchmarks are flexible, in that they allow researchers to explicitly
    modify the benchmark to identify particular task domains where particular
    hardware excels.  Furthermore, the benchmarks make use of a hybrid of
    real physical embodiment and a type of ``minimal'' simulation that has been
    shown to lead to robust real-world performance, allowing the same benchmark
    to be used by many researchers.

\tiny
 \keyFont{ \section{Keywords:} neuromorphic hardware, benchmarking, 
        minimal simulation, adaptive control, neural networks} %All article types: you may provide up to 8 keywords; at least 5 are mandatory.
\end{abstract}

\section{Introduction}

Neuromorphic hardware holds great promise for a wide variety of applications.
The combination of massively parallel computation and low power consumption
means that there is the potential to have complex algorithms running in
embedded processing situations, without being a significant drain on battery
life.  The question, of course, is to identify what sort of always-on or
interactive functionality is feasible with these devices.

To evaluate applications of this hardware, we need benchmark tasks.  These
tasks must allow us to compare across different instances of neuromorphic
hardware (and potentially across different algorithms implemented in said
hardware).  This allows us to quantitatively compare systems, letting 
researchers both measure the progress in the field and also directly
compare competing approaches.

In this paper, we focus on the development of \emph{embodied} benchmarks.
These are tasks where the output of the neuromorphic hardware \emph{influences
its own future input}.  This is in contrast to standard categorization or pattern
identification tasks, where the input is some fixed sequence, and the main
question is whether the hardware produces the correct output for each input
(or input pattern).

We believe embodied benchmarks should be of particular interest to
neuromorphic research, given that many of the applications of neuromorphic
hardware are in exactly this domain of embedded and interactive control
of robotic or other physical systems.  However, embodiment raises a number
of issues that complicate the development of such benchmarks.  Rather than
simply providing a data file of inputs and desired outputs, the benchmark
must either specify a full physical system for that embodiment, or it must 
provide software for a simulation of that system.  As we discuss below, either
approach is problematic, and addressing this difficulty is the primary goal of
this paper.


\section{Embodied Benchmarks}

An embodied benchmark task is one where the system we are studying has a 
two-way interaction with some sort of environment.  That is, the outputs from
the neuromorphic hardware are sent to the environment where they cause some
sort of effect, the results of which change the subsequent input.  For example,
the outputs might control the movement of a robot, which in turn affects the
sensory data received by the robot.

\subsection{Simulation versus Physical Instantiation}

To define such a benchmark, we need to be explicit about the embodiment.  If
it is a situation where a robot is to be controlled, we need to be explicit
about all the details of that robot.  What motors are there?  How are they
configured?  How strong are they?  What sensors are there?  Where are they
placed?  How accurate are they?  However, even if these questions are
answered, there is a fundamental problem in that \emph{other researchers
need access to that exact robot}.  If a benchmark is to be widely used, other
researchers developing their own hardware should be able to do their own 
testing.

Furthermore, using a physical robot also imposes significant practical
difficulties when performing extensive benchmark testing.  When testing, we
often want to run the same task over and over again, both for robustness and
to see the effects of varying parameters.  With a physical robot, this means
manually setting up the task, letting the test run, gathering the resulting
data, and then resetting the robot back to the initial state.  This means that
issues like battery life become problematic, and not just because there is a
limited amount of time available for testing.  As the battery level changes,
the performance of the robot itself can also change.  Futhermore, for any
rigorous testing of the benchmark, we will want to examine situations where
the system fails.  This means that some of the testing will involve parameter
settings that lead to poor behaviour, which could involve physical damage to
the robot itself.

However, \emph{not} using a real physical embodiment for testing is also
problematic.  First and foremost, without an actual real-world task, why
should anyone have any confidence that the performance on the benchmark is
reflective of the actual usefulness of the neuromorphic hardware?  It is
very widely known that \emph{simulations of robots} (or other physical systems) are
often much easier to control and better-behaved than the real thing.  The
field of robotics is filled with algorithms that work well ``in theory,'' but
fail when run on actual hardware.  We do not want a benchmark that falls into
this trap of failing to generalize to real situations.

Furthermore, neuromorphic hardware has another constraint that severely limits
the possibilities of simulation.  The normal approach when a simulation is
too simple to reflect reality is to add details to the simulation itself.
Incredibly finely detailed simulations can be created, filling in all of the
details needed.  However, those resulting simulations \emph{cannot be run
in real-time}.  This is a fundamental problem, in that most neuromorphic
hardware is about real-time interactions, and there is no way to slow down
the hardware to match the simulated environment.  This means that even if
we spent the considerable amount of research effort needed to define a
simulated environment for an embodied benchmark, that simulation could not
be run fast enough to interact with most of the hardware that we would want
to test.

\subsection{Minimal Simulation}

The above considerations seem to indicate that even though using real-world
phyiscal hardware for benchmarking is problematic, it is still better than
using simplistic simulations which may not generalize to real tasks.  However,
we believe there is an alternate approach known as \emph{Minimal Simulation}
(\citealt{Jakobi97evolutionaryrobotics}).

First, we note that the problem faced here is remarkably similar to a problem
faced by the evolutionary robotics community.  In evolutionary robotics, the
goal is to use genetic algorithms to \emph{evolve} systems that can control
robots to perform various tasks.  These tasks can be as simple as navigation
and obstacle avoidance, but have also included walking,
object identification, and visual tracking [refs???].

However, performing this evolution on real physical robots is problematic
for the same reasons that benchmarks on physical robots are problematic.  The
robots must be reset to the same state each time; they often involve
behaviour that can physically damage the robots; and they take a very long
time to run.  For this reason, attempts were made to evolve using simulated
robots.  However, the general finding was that algorithms that worked on the
simulated robots would not work when run on the real physical robots.  If
the simulations were improved, adding complex physical detail, then it was
possible to generalize to real behaviour; unfortunately, such complex
simulations would run slower than real-time [refs].

To address this problem, Nicholas Jacobi proposed the creation of
``minimal'' simulations.  These are simulations where there is variability
\emph{within the simulation itself}.  In other words, we make \emph{poor}
simulations, but ensure that the way in which it is poor is itself
variable.  We then make sure that the controllers work across that whole
range of variability.  ``Instead of trying to eliminate the differences between
simulation and reality, they are acknowledged, and mechanisms are put in place
to prevent evolving controllers from relying on them.'' [ref: jacobi thesis].

With this approach, it was possible to build minimal simulations that would
run faster than real-time and yet also be complex enough that if a system
could successfully control the simulation, it was also likely to successfully
control a real robot.  To achieve this, the simulations are made to be
unreliable in almost every respect.  For example, for a simulation of a simple
motor it would still be the case that if power is applied it would generally
try to spin, but the exact amount of torque, the amount of sensory noise,
the amount of time needed, the amount of static and dynamic friction, and so
on would all be randomly chosen.  A successful controller would have to deal
with this wide range of variability, and if it could handle that variability
then we would have reason to believe it could also handle the real system.

It is also worth noting that a minimal simulation \emph{only has to be a
good simulation for successful behaviour}.  That is, ``if we are evolving 
corridor following behaviour, the dynamics of the simulation might
differ wildly from those of reality if the controller hits a wall or
goes round in circles, but this does
not matter, since the controllers we are interested in transferring
across the reality gap will neither
hit walls nor go round in circles.''  If the controller is poor, we do not 
need the simulation to be at all accurate in exactly \emph{how} that poor
behaviour is manifest.  We do not need an exact detailed physics model of
the collision between a robot and a wall, or a detailed model of what happens
to a robot arm when it starts oscillating wildly due to a poor control signal.
All we need is for the simulation to be just good enough to indicate that
things have gone wrong, and thus give a low score to that controller.  This
means that, for example, in a minimal simulation of an eight-legged walking robot,
it is not necessary to have a physics simulation that correctly models what
happens when two legs collide with each other.  Rather, if legs collide
with each other, that is an indication that the walking behaviour is very poor,
so as long as that result is indicated we can greatly simplify the simulation
by not including all the details necessary to deal with these sorts of physical
interactions.

Given the success of this approach for evolutionary robotics, we propose
that it can be directly used for embodied neuromorphic benchmarks.  To do
this, we create software simulations for each benchmark task.  These simulations
must be fast enough to run in realtime (so that they can be controlled by
real neuromorphic hardware), and they must be extremely variable.  Each time
the simulation is run, different parameters will be chosed for this variability
(so one run might have a large degree of sensor noise while the next run has
none at all; one run might have more delay in the motor response and another
might have less power available).  Being successful at the benchmark means
being successful across all this variability.  

The result should be a benchmark that can be run by any researcher.  The fact
that it is a simulation means that source code can be shared, and that no
specialized hardware is needed.  Furthermore, the variability in the simulation
itself can be controlled, and this can help give a rich characterization of
the benchmarked hardware.  For example, some hardware might only work with
small amounts of sensor noise, or other hardware might only work when there
is significant delay in the motor response.  This flexibility in parameters
in the benchmark allows researchers to explicitly characterize that particular
siuations where their hardware excells.  

\subsection{Cheap Robotics}

All of that said, we also need real physical embodiment as part of any
benchmark of this kind.  While we claim that reliable minimal simulations
are relatively easy to create, they are not a widely-used technique, and the
real proof of behaviour is always in the real world.  For this reason, we
also argue that every benchmark based on a minimal simulation should also
have an easy-to-construct physical analog.  However, this physical version
is not to be used as the primary benchmark.  Rather, this is meant as a
double-check that the system does actually behave as expected.  

For this physical aspect of the benchmark, we recommend cheap, widely-available
components.  This allows a greater chance for other researchers to have
access to the same (or similar) hardware.  For the particular example 
benchmark described in the next section, we use the Lego Mindstorms EV3 kit,
a simple robotics platform available at most toy stores.

It is important to note that there is actually a theoretical advantage to
using cheap robotics hardware for benchmarking, in addition to the practical
advantages.  In particular, we \emph{don't want benchmarks that rely on
particular high-speed, high-accuracy devices}.  The purpose of the benchmark
is not to indicate how well this neuromorphic hardware works to control this
one particular robot in this task.  Rather, the purpose of a benchmark is to
characterise how well this neuromorphic hardware works on this task 
\emph{in general}.  The variability in the minimal simulation means that it
should be able to function across a wide variety of physical embodiments, and
so if we are to choose one particular physical embodiment to test in the
real world, then we should choose one that is not extremely high-precision.
For this reason, we believe using cheap Lego robot is actually more useful
for benchmarking than an expensive high-precision robot.\footnote{Of course,
for more complex benchmark tasks we may need sensory and motor capabiities
that are beyond that of a simple Lego robot.}

\section{A Benchmark: Adaptive Motor Control}

To demonstrate this approach to creating embodied neuromorphic benchmarks,
we now consider a basic control task.  Suppose we have a system with a number
of joints $q$ and we want to send an output $u$ to those motors such that
the joints move to a particular desired position $q_d$.  Our only output is
the signal $u$ (one for each motor) and our only input is the current position
of each motor $q$.

The simplest controller for such a situation is a P (proportional) controller,
where $u=K_p(q_d - q)$.  This is often supplemented with a D (derivative) term,
which helps to slow the system down as it approaches the desired position,
thus avoiding overshooting and oscillation ($u=K_p(q_d - q) + K_d(\dot{q_d} - \dot{q})$),
leading to the standard PD controller.  Both $K_p$ and $K_d$ are constants that can be
tuned to particular situations.

However, this controller has difficulty in the presence of significant external
forces.  For example, consider a single motor controlling the angle of a single
arm.  If the arm is held out to the side, gravity acting on the mass of the
arm itself will pull the arm downward.  Thus to hold the arm still at a
particular $q_d$ will require the controller to apply a force to counteract
gravity.  Since the PD controller always produces an output $u=0$ when $q=q_d$,
it cannot compensate for this, and so the arm will stay stationary at some
angle below the desired angle. [TODO: add diagram]

The standard solution to this problem is to add an I (integral) term ($K_i \int{(q_d-q) dt}$)
to the controller, making it a PID controller.  The idea here is that as the difference between where it is and
where we want it to be accumulates over time, the $K_i$ term gradually increases
how much extra force is being applied until it is large enough to counteract
the external force of gravity (or whatever other external forces are present).
However, this approach has great difficulty when $q_d$ changes, since the
external force due to gravity changes depending on the position of the arm $q$.
The controller ends up having to use the accumulated integral to ``relearn'' the
correct amount of extra force needed every time $q_d$ changes.

In some robotics applications, the solution to this problem is to mathematically
analyze the geometry and mass of the system and compute exactly how much extra
force is needed.  In this particular case, the answer is straight-forward,
in that the extra torque due to gravity is $\tau=mg {l \over 2} sin(q)$, where
$m$ is the mass of the arm, $l$ is the length, and $g$ is $9.8m/s^2$.  If the
force applied by the motor is linear in $u$, then we could simply compute this
value and add it to our controller's output.  However, this assumes a perfectly
even distribution of weight in the arm, ignores momentum and other forces, and 
gets much more complex as more joints are added.  Furthermore, if this
initial computation is slightly off, or if details of the system change,
there is no way to adjust this compensation.

Fortunately, there is an adaptive solution to this problem, and it is one that
fits well with neuromorphic hardware.  \cite{Slotine1987} shows that if you
express the influence of these other external forces as $\tau=Y(q) \omega$
(where $Y(q)$ is a fixed set of functions of $q$, such as $sin(q)$, and $\omega$
is a vector of scalar weights, one for each function in $Y$), then you can learn
to compensate for these external forces by using the learning rule $\Delta \omega = \alpha Y(q) u$,
where $u$ is the basic PD control signal.

Importantly, as pointed out by \cite{SannerSlotine1992} and \cite{Lewis1996}, rather than making explicit
assumptions about the exact functions that should be in $Y(q)$, we can use
a neural network approach where each neuron is a different function of $q$.  
As long as there is enough hetereogenetity (i.e. as long as the neural activity
forms a basis space that spans the desired space of functions), then the
learning rule will continue to work.  This approach has been used in the
Recurrent Error-driven Adaptive Control Hierarchy model of human motor control
(\citealt{dewolf2014}), and for quadcopter control [TODO: any citations for this?].

This then suggests an explicit neuromorphic benchmark.  The input to the
neuromorphic hardware is $q$, the system state.  This input is fed to each
neuron such that each neuron produces some output behaviour that is based on
this input.  Since $q$ will be multi-dimensional (if there is more than one
joint), we may give each neuron a random weighting of each $q$ value ($J_i=e_i \cdot q$, 
where $J_i$ is the input to neuron $i$, and $e_i$ is a randomly chosen vector\footnote{$e$ could also be chosen so as to regularly span the space of possibilities}).  Given this input, the neurons will produce some output $A$.
we now form a weighted sum of these outputs $Ad$, where $d$ is a matrix (number of neurons by number of elements in $q$)
that is initially all zeros.

To use this controller, we add its output to that of the standard PD controller.
That is, the standard controller has $u=K_p(q_d - q) + K_d(\dot{q_d} - \dot{q})$,
and so our actual output to the motor is $u + Ad$.  We then apply a learning
rule on $d$ such that $\Delta d = \alpha A \times u$.

Notice that we can think of this system as a three-layer neural
network, where the input and output layers are linear.  The first layer is $q$, the input state, one value for each joint.  The ``hidden'' layer
is $A$, the activity of a large number of neurons.  The output layer again has
one value per joint, and is the extra added signal to apply to the motors, $Ad$.
Given that this is such a canonical example of the use of neural networks, we
hope that the majority of neuromorphic hardware is flexible enough to implement
exactly this model.


\subsection{Online and offline learning}

The one major step here that does not exist in a lot of neuromorphic hardware is
the ability to update the weights $d$.  For hardware that does have a built-in
learning rule, this rule is at least of a very common form, where the weight
update from a neuron is proportional to the activity of that neuron and an
external error signal.  This makes it an instance of the ubiquitous delta rule,
and hopefully supported by the hardware.

However, if the neuromorphic hardware being benchmarked does not have the
ability to update weights online using a learning rule of this form, then there
are two solutions.  First, the multiplication by $d$ could be done on the
output from the neuromorphic hardware.  There has to be some system to take
the neural output from the hardware and send it to the motor (or the simulation
of the motor).  Instead of outputing the result of $Ad$, the hardware could
output $A$ (the activity of all the neurons), and the interface to the motor
could be responsible for doing the multiplication by $d$ and updating $d$
according to the learning rule.

Alternatively, we can use offline learning.  That is, rather than updating
the weights $d$ all the time, we simply record $A$ and $u$, and then after
a period of time stop the controller, compute the sum total of the changes
to $d$, load the new value of $d$ onto the neuromorphic hardware, and then
start the controller again.

\subsection{Minimal Simulation for Adaptive Control}

Now that we have defined the task, we can use the principles of Minimal
Simulation to construct a flexible and variable simulated environment for
testing adaptive control.  The idea is to make a bare-bones simulation of
the system being controlled, with built-in variability.  If the neuromorphic
controller works well across this variability, then it is likely to work well
outside of simulation as well.

The basic system variable is a vector of length $N$ holding the joint angles $q$.  
Each joint has a velocity $v$.  The force applied by each motor is 
related to the signal sent to the motor $u$, but will generally have some
maximum value $T$, so we use $tanh(u)T$. For friction, we 
simply scale the velocity by some factor $F$ every time step.  
This results in the simplistic simulation as follows:

\begin{align}
    \Delta v & =  -v F + tanh(u)T \\
    \Delta q & = v
\end{align}

On top of this, we need to add the external force for which the adaptive system
is trying to compensate.  Rather than deciding on what specific external
force, we \emph{randomly generate} an equation controlling this force.  We start with a small set of
smooth functions on $X$ which are often found in dynamics equations ($x$, $x^2$,
$sin(x)$).  We then generate an external force of $Q (\zeta X(\beta q + \gamma) + \eta)$
where $\zeta$, $\beta$, $\gamma$, and $\eta$ are all randomly chosen fron $N(0,1)$
and $Q$ is a scaling factor to control how strong this external force is.  The
result is added to Equation 1.

Finally, we add random noise, delay, and filtering to both the input and the
output of the system.  For noise, we add $N(0,\sigma_u)$ to the control signal
$u$ and $N(0,\sigma_q)$ to the $q$ value reported back to the controller.  We
also use a low-pass filter to smooth both values (with time constants $\tau_u$
and $\tau_q$) after this noise is added.  Finally, both $q$ and $u$ are delayed
by an amount of time $t_q$ and $t_u$.

The result is not meant to be a simulation of a particular physical embodiment.
Rather, the variability in this simulation is meant to be extremely fast
to simulation, and to make it hard to ``cheat''
to control it.  In other words, if a controller manages to be able to control
the various randomly created minimal simulations of embodiment that are
generated with this approach, then we have reason to believe that it will also
be successful at controlling real embodiments.  With this simulation and
modern computers, we can easily simulate in real-time systems with dozens of joints and
highly complex interactions between them, and measure how well an adaptive
controller deals with these situations.

\subsection{Cheap Robotics for Adaptive Control}

Of course, we need physical confirmation that the minimal simulation defined
in the previous section does a reasonable job.  For this, we define a
particular easy-to-build example of a system that can be usefully controlled
by this adaptive method.  For this, we use the Lego Mindstorms EV3 robot kit.

The physical body is depicted in [TODO: add figure].  It consists of a single
motor, mounted in such a way that there is a significant amount of weight on
the arm itself (the weight of the motor itself).  Multiple motors can be
added, and other configurations can be considered and should also be suitable
for testing control, but here we just consider the basic case.

To interface to the physical hardware, we installed the \texttt{ev3dev} operating
system (\texttt{http://ev3dev.org}), a Debian-based Linux system specifically developed
for the EV3.  We then installed and ran the \texttt{ev3\_link} program from \texttt{ev3dev-c}
(\texttt{https://github.com/in4lio/ev3dev-c}).  This allows the EV3 to listen for
UDP commands that tell it to set motor values and read sensor values.
Communication with a PC was over a USB link (although the system also
supports WiFi communication).  With constant communication, the system is
able to adjust the power sent to the motors $u$ and give position feedback
$q$ from those motors at a rate of around 300Hz.

Figure [???] shows the system being controlled by a simple PD controller.
The desired joint angle $q_d$ is varied randomly over time.  Notice that the
system consistently undershoots and overshoots the desired angle, due to
its inability to compensate for gravity.

In Figure [???], we show the result of using the adaptive controller algorithm.
The system quickly improves on the standard PD controller, learning to adjust
its control signal $u$ to compensate for the unknown external forces.

\section{Benchmark Analysis}

- what do we measure?   performance: rmse and delay.  Also, power consumption?  cost?

- for given hardware, how many neurons can be done in realtime?  that's how many should be used for comparison

- show how these things change as we change parameters in the simulation

- delays

- number of motors $N$

- magnitude of nonlinearity $Q$

CPU, GPU, and SpiNNaker?

\section{Other Benchmarks}

(sketch out how this approach could be applied to other task)

Adaptive Jacobian

Driving down a corridor, avoiding obstacles

Classical conditioning

Operant conditioning 


\section*{Disclosure/Conflict-of-Interest Statement}

The authors declare that the research was conducted in the absence of any commercial or financial relationships that could be construed as a potential conflict of interest.

\section*{Acknowledgments}
Text Text Text Text Text Text  Text Text Text Text Text Text Text Text  Text Text Text Text Text Text Text Text Text  Text Text Text. Text Text Text Text Text Text  Text Text Text Text Text Text Text Text  Text Text Text Text Text Text Text Text Text  Text Text Text. 

\textit{Funding\textcolon} Text Text Text Text Text Text  Text Text.

\bibliographystyle{frontiersinSCNS_ENG_HUMS} % for Science, Engineering and Humanities and Social Sciences articles, for Humanities and Social Sciences articles please include page numbers in the in-text citations
%\bibliographystyle{frontiersinHLTH&FPHY} % for Health and Physics articles
\bibliography{paper}

%%% Upload the *bib file along with the *tex file and PDF on submission if the bibliography is not in the main *tex file

\section*{Figures}

%%% Use this if adding the figures directly in the mansucript, if so, please remember to also upload the files when submitting your article
%%% There is no need for adding the file termination, as long as you indicate where the file is saved. In the examples below the files (logo1.jpg and logo2.eps) are in the Frontiers LaTeX folder
%%% If using *.tif files convert them to .jpg or .png

\begin{figure}[h!]
\begin{center}
\includegraphics[width=10cm]{logo1}% This is a *.jpg file
\end{center}
 \textbf{\refstepcounter{figure}\label{fig:01} Figure \arabic{figure}.}{ Enter the caption for your figure here.  Repeat as  necessary for each of your figures }
\end{figure}

%\begin{figure}
%\begin{center}
%\includegraphics[width=10cm]{logo2}% This is an *.eps file
%\end{center}
%\textbf{\refstepcounter{figure}\label{fig:02} Figure \arabic{figure}.}{ Enter the caption for your figure here.  Repeat as  necessary for each of your figures }
%\end{figure}

%%% If you don't add the figures in the LaTeX files, please upload them when submitting the article.

%%% Frontiers will add the figures at the end of the provisional pdf automatically %%%

%%% The use of LaTeX coding to draw Diagrams/Figures/Structures should be avoided. They should be external callouts including graphics.

\end{document}
